\newgeometry{left=1.25in,right=1.25in,top=1in,bottom=1in}
\phantomsection
\addcontentsline{toc}{chapter}{一、文献阅读卡}
\label{chap:literature-cards}
\setlength{\parindent}{0pt}
\setlength{\parskip}{0pt}
\setstretch{1.5}

\begin{center}
{\songti\bfseries\fontsize{18.5pt}{27.75pt}\selectfont 华中科技大学设计学院硕士研究生\par}
{\songti\bfseries\fontsize{18.5pt}{27.75pt}\selectfont 文献阅读卡\par}
\end{center}
\vspace{0.25em}
\noindent\underline{\makebox[\textwidth][l]{\textbf{学号：M202570493\hfill 姓名：陈思静\hfill 序号：01}}}\par
\vspace{0.35em}
\noindent\textbf{名称}\ Results-Actionability Gap: Understanding How Practitioners Evaluate LLM Products in the Wild\par
\noindent\textbf{作者}\ Willem van der Maden; Malak Sadek; Ziang Xiao; Aske Mottelson; Q. Vera Liao; Jichen Zhu\quad 译者\par
\noindent\textbf{出处}\ 【美国】Proceedings of the 2026 CHI Conference on Human Factors in Computing Systems，2026：1–17\quad 页\par
\vspace{0.35em}
\noindent\textbf{文献阅读心得（每篇不少于150字）：}\par
\noindent\textbf{研究背景}\ 大语言模型（large language model，LLM）产品的评价对象不是孤立模型，而是包含提示词、检索、界面与业务流程的完整系统。传统基准分数很难直接回答产品团队下一步应改什么，实践者因此在不确定性、资源限制和多方目标之间组织评价工作。\par
\noindent\textbf{研究问题}\ 研究问题一（RQ1）：生产环境中的大语言模型产品目前采用哪些评价实践？\par
研究问题二（RQ2）：实践者认为评价大语言模型产品时最主要的困难是什么？\par
\noindent\textbf{研究方法}\ 作者对19名跨行业大语言模型产品实践者开展半结构化访谈，并进行主题分析。研究覆盖评价执行、评价方案设计及组织层面的协调工作，分析对象是面向真实用户的完整产品。\par
\noindent\textbf{主要发现}\ 研究归纳出10类评价活动和5类挑战。17名参与者经历了“结果—可行动性鸿沟”：团队能收集反馈或分数，却难以判断失败来自提示、检索、模型还是其他部件。人工判断与“直觉检查”（vibe check）被视为对大语言模型特性的适应，而不只是方法不成熟。\par
\noindent\textbf{贡献与意义}\ 论文提出“规范化历程”（formalization journey）和“结果—可行动性鸿沟”（results-actionability gap），把评价后的定位、协商与改进纳入评价研究，并从较成熟团队中总结了可实施的组织与流程策略。\par
\clearpage
\begin{center}
{\songti\bfseries\fontsize{18.5pt}{27.75pt}\selectfont 华中科技大学设计学院硕士研究生\par}
{\songti\bfseries\fontsize{18.5pt}{27.75pt}\selectfont 文献阅读卡\par}
\end{center}
\vspace{0.25em}
\noindent\underline{\makebox[\textwidth][l]{\textbf{学号：M202570493\hfill 姓名：陈思静\hfill 序号：02}}}\par
\vspace{0.35em}
\noindent\textbf{名称}\ Interactive Debugging and Steering of Multi-Agent AI Systems\par
\noindent\textbf{作者}\ Will Epperson; Gagan Bansal; Victor Dibia; Adam Fourney; Jack Gerrits; Erkang Zhu; Saleema Amershi\quad 译者\par
\noindent\textbf{出处}\ 【美国】Proceedings of the 2025 CHI Conference on Human Factors in Computing Systems，2025：156:1–156:15\quad 页\par
\vspace{0.35em}
\noindent\textbf{文献阅读心得（每篇不少于150字）：}\par
\noindent\textbf{研究背景}\ 多智能体（multi-agent）系统会产生数十甚至上百条消息，错误可能在早期计划、角色分工或工具调用中出现，并沿后续协作传播。现有开发工具多支持事后阅读日志，缺少暂停、回退和试改轨迹的交互能力。\par
\noindent\textbf{研究问题}\ 研究问题一（RQ1）：开发者在构建和调试多智能体工作流时面临哪些核心困难？\par
研究问题二（RQ2）：消息级浏览、编辑与重置如何支持故障诊断和智能体引导（agent steering）？\par
\noindent\textbf{研究方法}\ 研究先访谈5名智能体开发者，据此设计交互调试工具AGDebugger；随后开展含14名参与者的两阶段用户研究。第一阶段比较AGDebugger与基线日志界面的错误识别，第二阶段让参与者编辑历史消息并重新运行。\par
\noindent\textbf{主要发现}\ 参与者通常花约10分钟阅读日志后才开始修改。5/6的第一阶段参与者偏好AGDebugger；第二阶段8人共做24次编辑，其中14次是补充具体指令，另有简化指令和修改计划。重置与编辑功能评分最高，但只有2/8得到完全正确答案。\par
\noindent\textbf{贡献与意义}\ AGDebugger把历史概览、逐条消息检查、回退、编辑与继续执行放进同一界面，使开发者能够用反事实式试验观察某次干预是否改变后续轨迹。\par
\clearpage
\begin{center}
{\songti\bfseries\fontsize{18.5pt}{27.75pt}\selectfont 华中科技大学设计学院硕士研究生\par}
{\songti\bfseries\fontsize{18.5pt}{27.75pt}\selectfont 文献阅读卡\par}
\end{center}
\vspace{0.25em}
\noindent\underline{\makebox[\textwidth][l]{\textbf{学号：M202570493\hfill 姓名：陈思静\hfill 序号：03}}}\par
\vspace{0.35em}
\noindent\textbf{名称}\ Which Agent Causes Task Failures and When? On Automated Failure Attribution of LLM Multi-Agent Systems\par
\noindent\textbf{作者}\ Shaokun Zhang; Ming Yin; Jieyu Zhang; Jiale Liu; Zhiguang Han; Jingyang Zhang; Beibin Li; Chi Wang; Huazheng Wang; Yiran Chen; Qingyun Wu\quad 译者\par
\noindent\textbf{出处}\ 【美国】Proceedings of Machine Learning Research，2025，267：76583–76599\quad 页\par
\vspace{0.35em}
\noindent\textbf{文献阅读心得（每篇不少于150字）：}\par
\noindent\textbf{研究背景}\ 多智能体任务失败时，最后输出错误的智能体不一定是决定结果的源头。有效调试需要同时定位责任智能体和关键步骤，但长轨迹、角色交互和错误传播使人工归因成本很高。\par
\noindent\textbf{研究问题}\ 研究问题一（RQ1）：能否自动识别导致任务失败的智能体，即回答“谁”造成失败？\par
研究问题二（RQ2）：能否定位该智能体产生决定性错误的具体步骤，即回答“何时”出错？\par
\noindent\textbf{研究方法}\ 作者将修正某个“智能体—时间”对后能否把失败变为成功作为决定性错误的反事实定义，并以最早的决定性错误为标注目标。Who\&When数据集包含127个多智能体系统、184个失败归因任务，比较一次性分析（all-at-once）、二分搜索（binary search）和逐步分析（step-by-step）三类方法。\par
\noindent\textbf{主要发现}\ 最佳方法的智能体级准确率为53.5\%，步骤级准确率仅14.2\%，部分方法低于随机基线；包括先进推理模型在内的方法仍未达到实用水平。自动判断“谁”已有有限信号，而精确判断“何时”更困难。\par
\noindent\textbf{贡献与意义}\ 论文首次系统形式化多智能体自动故障归因，提供同时含责任智能体、决定性步骤和自然语言原因的细粒度数据集，并把智能体级与步骤级能力分开评价。\par
\clearpage
\begin{center}
{\songti\bfseries\fontsize{18.5pt}{27.75pt}\selectfont 华中科技大学设计学院硕士研究生\par}
{\songti\bfseries\fontsize{18.5pt}{27.75pt}\selectfont 文献阅读卡\par}
\end{center}
\vspace{0.25em}
\noindent\underline{\makebox[\textwidth][l]{\textbf{学号：M202570493\hfill 姓名：陈思静\hfill 序号：04}}}\par
\vspace{0.35em}
\noindent\textbf{名称}\ Finding AI's Faults with AAR/AI: An Empirical Study\par
\noindent\textbf{作者}\ Roli Khanna; Jonathan Dodge; Andrew Anderson; Rupika Dikkala; Jed Irvine; Zeyad Shureih; Kin-Ho Lam; Caleb R. Matthews; Zhengxian Lin; Minsuk Kahng; Alan Fern; Margaret Burnett\quad 译者\par
\noindent\textbf{出处}\ 【美国】ACM Transactions on Interactive Intelligent Systems，2022，12(1)：1:1–1:33\quad 页\par
\vspace{0.35em}
\noindent\textbf{文献阅读心得（每篇不少于150字）：}\par
\noindent\textbf{研究背景}\ 总体成败或平均性能无法告诉领域用户人工智能究竟在哪些情境下不可靠。要决定何时依赖智能体，用户需要把问题定位到具体决策和内部预测，而不只是给出笼统的通过或失败判断。\par
\noindent\textbf{研究问题}\ 研究问题一（RQ1）：人工智能行动后复盘流程（After-Action Review for AI，AAR/AI）能否帮助具备领域知识的用户定位人工智能故障？\par
研究问题二（RQ2）：故障类型不同，用户发现故障的效果是否也不同？\par
\noindent\textbf{研究方法}\ 65名有即时战略游戏经验但无人工智能或机器学习背景的参与者被随机分到复盘组和对照组，分析强化学习智能体完成《星际争霸II》拉锯战任务的过程。研究用10个经验证的自然故障实例计算问题报告的召回率和精确率。\par
\noindent\textbf{主要发现}\ AAR/AI组召回率为0.233，对照组为0.080；精确率分别为0.179和0.116，差异均显著。AAR/AI参与者在每类故障上表现更好，发现任一特定故障的可能性平均接近对照组的6倍。\par
\noindent\textbf{贡献与意义}\ 论文用可核对的故障真值、召回率和精确率评价人类诊断，而非只测满意度；结构化复盘和标签活动还可能帮助用户从单个实例上升到故障类别。\par
\clearpage
\begin{center}
{\songti\bfseries\fontsize{18.5pt}{27.75pt}\selectfont 华中科技大学设计学院硕士研究生\par}
{\songti\bfseries\fontsize{18.5pt}{27.75pt}\selectfont 文献阅读卡\par}
\end{center}
\vspace{0.25em}
\noindent\underline{\makebox[\textwidth][l]{\textbf{学号：M202570493\hfill 姓名：陈思静\hfill 序号：05}}}\par
\vspace{0.35em}
\noindent\textbf{名称}\ Debugging Reinvented: Asking and Answering Why and Why Not Questions about Program Behavior\par
\noindent\textbf{作者}\ Andrew J. Ko; Brad A. Myers\quad 译者\par
\noindent\textbf{出处}\ 【美国】Proceedings of the 30th International Conference on Software Engineering，2008：301–310\quad 页\par
\vspace{0.35em}
\noindent\textbf{文献阅读心得（每篇不少于150字）：}\par
\noindent\textbf{研究背景}\ 开发者观察到异常输出后，传统调试器仍要求他们先猜测相关代码，再用断点和单步执行验证。由于初始猜测经常错误，问题表述与工具操作之间的转换会产生大量无关搜索。\par
\noindent\textbf{研究问题}\ 研究问题一（RQ1）：如何从程序代码和一次执行中自动提取用户可能提出的“为什么”与“为什么没有”（why and why-not）问题？\par
研究问题二（RQ2）：如何用静态、动态分析回答这些问题，并把答案映射到相关运行事件和代码？\par
\noindent\textbf{研究方法}\ 论文给出Java Whyline的系统设计与算法：记录完整执行历史，从文本和图形输出生成问题，结合动态切片、静态切片、调用图、可达性分析和值来源分析构造因果解释，并讨论性能与覆盖限制。\par
\noindent\textbf{主要发现}\ “为什么发生”可沿已执行依赖回溯；“为什么没有发生”必须分析潜在定义、未满足条件与未到达路径，计算和问题空间更复杂。原型证明可以从可见输出进入相关代码，而无需先猜故障位置。\par
\noindent\textbf{贡献与意义}\ 文章把调试问题作为一等交互对象，并提出面向大型Java程序自动生成、筛选和回答“为什么”与“为什么没有”问题的完整技术架构。\par
\clearpage
\begin{center}
{\songti\bfseries\fontsize{18.5pt}{27.75pt}\selectfont 华中科技大学设计学院硕士研究生\par}
{\songti\bfseries\fontsize{18.5pt}{27.75pt}\selectfont 文献阅读卡\par}
\end{center}
\vspace{0.25em}
\noindent\underline{\makebox[\textwidth][l]{\textbf{学号：M202570493\hfill 姓名：陈思静\hfill 序号：06}}}\par
\vspace{0.35em}
\noindent\textbf{名称}\ Designing the Whyline: A Debugging Interface for Asking Questions about Program Behavior\par
\noindent\textbf{作者}\ Andrew J. Ko; Brad A. Myers\quad 译者\par
\noindent\textbf{出处}\ 【美国】Proceedings of the 2004 ACM Conference on Human Factors in Computing Systems，2004：151–158\quad 页\par
\vspace{0.35em}
\noindent\textbf{文献阅读心得（每篇不少于150字）：}\par
\noindent\textbf{研究背景}\ 调试本质上从对运行行为的疑问开始，但断点、单步和打印语句迫使程序员先定位代码。Whyline尝试让用户直接针对输出提问，再由系统呈现与该问题有关的运行时因果链。\par
\noindent\textbf{研究问题}\ 研究问题一（RQ1）：用户是否认为基于“为什么”与“为什么没有”提问的Whyline有用？\par
研究问题二（RQ2）：Whyline能否缩短调试时间并帮助程序员完成更多任务？\par
\noindent\textbf{研究方法}\ 作者在Alice编程环境中实现Whyline，并对既有无Whyline研究和新增有Whyline研究中的相同场景进行比较。无工具组4人、有工具组5人，在90分钟内完成6项Pac-Man编程任务并出声思考。\par
\noindent\textbf{主要发现}\ 在6个可比调试场景中，Whyline把调试时间平均缩短7.8倍；有工具组平均完成3.20项任务，无工具组为2.25项，提升约40\%。19/24次提问中工具被认为有用，“为什么没有”问题又明显多于“为什么”问题。\par
\noindent\textbf{贡献与意义}\ Whyline把可见程序输出转成问题菜单，并用时间游标、因果箭头和可逆历史把现象、运行事件与代码连接起来，改变了从代码位置开始的调试顺序。\par
\clearpage
\begin{center}
{\songti\bfseries\fontsize{18.5pt}{27.75pt}\selectfont 华中科技大学设计学院硕士研究生\par}
{\songti\bfseries\fontsize{18.5pt}{27.75pt}\selectfont 文献阅读卡\par}
\end{center}
\vspace{0.25em}
\noindent\underline{\makebox[\textwidth][l]{\textbf{学号：M202570493\hfill 姓名：陈思静\hfill 序号：07}}}\par
\vspace{0.35em}
\noindent\textbf{名称}\ Finding Causes of Program Output with the Java Whyline\par
\noindent\textbf{作者}\ Andrew J. Ko; Brad A. Myers\quad 译者\par
\noindent\textbf{出处}\ 【美国】Proceedings of the 27th ACM Conference on Human Factors in Computing Systems，2009：1569–1578\quad 页\par
\vspace{0.35em}
\noindent\textbf{文献阅读心得（每篇不少于150字）：}\par
\noindent\textbf{研究背景}\ 大型Java程序中，断点调试要求开发者先选对代码位置，而用户通常更容易识别界面中出现或缺失的输出。Java Whyline把“从结果反查原因”的思路扩展到真实开源项目。\par
\noindent\textbf{研究问题}\ 研究问题一（RQ1）：Java Whyline是否比断点调试器更快、更成功地帮助开发者定位真实缺陷？\par
研究问题二（RQ2）：两种工具如何改变开发者浏览文件、使用问题和接近关键函数的路径？\par
\noindent\textbf{研究方法}\ 研究采用组间实验，两组各10名软件工程硕士生，分别使用Whyline或具有相同外观的断点调试器，诊断ArgoUML的两个真实缺陷。每项任务限时30分钟，并记录成功率、时间、文件访问和导航距离。\par
\noindent\textbf{主要发现}\ 任务1中Whyline组10/10成功，对照组3/10，且前者约快2倍；任务2中分别为4/10和0/10。总体上Whyline用户约有3倍成功率，并更快接近关键代码；但工具更善于找到故障方法，不一定能解释方法为何错误。\par
\noindent\textbf{贡献与意义}\ 研究把输出选择、执行录像、依赖可视化、书签与问题导航集成为可实证比较的调试界面，并使用真实大型项目而非人工小程序。\par
\clearpage
\begin{center}
{\songti\bfseries\fontsize{18.5pt}{27.75pt}\selectfont 华中科技大学设计学院硕士研究生\par}
{\songti\bfseries\fontsize{18.5pt}{27.75pt}\selectfont 文献阅读卡\par}
\end{center}
\vspace{0.25em}
\noindent\underline{\makebox[\textwidth][l]{\textbf{学号：M202570493\hfill 姓名：陈思静\hfill 序号：08}}}\par
\vspace{0.35em}
\noindent\textbf{名称}\ How Programmers Debug, Revisited: An Information Foraging Theory Perspective\par
\noindent\textbf{作者}\ Joseph Lawrance; Christopher Bogart; Margaret Burnett; Rachel Bellamy; Kyle Rector; Scott D. Fleming\quad 译者\par
\noindent\textbf{出处}\ 【美国】IEEE Transactions on Software Engineering，2013，39(2)：197–215\quad 页\par
\vspace{0.35em}
\noindent\textbf{文献阅读心得（每篇不少于150字）：}\par
\noindent\textbf{研究背景}\ 调试大型代码库时，程序员不断判断下一条线索是否值得追踪。仅用“提出并验证假设”描述这一过程，很难解释他们在文件、方法和调用关系之间的实际导航，也难为工具设计提供直接建议。\par
\noindent\textbf{研究问题}\ 研究问题一（RQ1）：专业程序员调试时的语言和导航是否主要体现信息气味追踪，而非显式假设？\par
研究问题二（RQ2）：信息气味与代码拓扑模型能否预测程序员下一步会访问哪里？\par
\noindent\textbf{研究方法}\ 作者提出基于信息觅食理论的调试模型，并观察10名专业程序员调试真实开源程序，分析出声思考与导航日志；随后构建可执行模型，与不考虑信息气味的替代模型比较预测能力。\par
\noindent\textbf{主要发现}\ 参与者的言语更多涉及寻找和比较线索，而不是陈述完整假设。错误信息、标识符、文件名和调用关系形成信息气味；纳入气味的模型更准确地预测导航，代码拓扑则决定线索间移动成本。\par
\noindent\textbf{贡献与意义}\ 论文以信息气味（scent）和代码拓扑（topology）两个可观测、可计算的构念解释调试导航，减少对难以直接测量的心理假设的依赖，并把理论转成可执行预测模型。\par
\clearpage
\begin{center}
{\songti\bfseries\fontsize{18.5pt}{27.75pt}\selectfont 华中科技大学设计学院硕士研究生\par}
{\songti\bfseries\fontsize{18.5pt}{27.75pt}\selectfont 文献阅读卡\par}
\end{center}
\vspace{0.25em}
\noindent\underline{\makebox[\textwidth][l]{\textbf{学号：M202570493\hfill 姓名：陈思静\hfill 序号：09}}}\par
\vspace{0.35em}
\noindent\textbf{名称}\ Mental Procedures in Real-Life Tasks: A Case Study of Electronic Trouble Shooting\par
\noindent\textbf{作者}\ Jens Rasmussen; A. Jensen\quad 译者\par
\noindent\textbf{出处}\ 【英国】Ergonomics，1974，17(3)：293–307\quad 页\par
\vspace{0.35em}
\noindent\textbf{文献阅读心得（每篇不少于150字）：}\par
\noindent\textbf{研究背景}\ 熟练维修人员面对电子设备故障时，不会穷举所有测量，而是在系统、功能单元和部件组成的层级结构中搜索。真实工作还要求兼顾诊断时间、认知负荷和测量成本。\par
\noindent\textbf{研究问题}\ 研究问题一（RQ1）：熟练电子维修人员如何把设备表示为层级系统并组织故障搜索？\par
研究问题二（RQ2）：不同搜索例程在观察数量、知识深度和心理加工复杂度上有何差异？\par
\noindent\textbf{研究方法}\ 研究在维修人员的正常工作环境中收集口述协议，对真实电子排故过程进行细致编码，分析症状判断、层级切换、测量选择及反复出现的搜索例程。\par
\noindent\textbf{主要发现}\ 排故可分解为多种循环搜索例程，分别依赖不同深度的内部机理知识。维修人员偏好快速连续的好/坏判断，即使某些观察在信息论上冗余；若目标是减少时间和心理负担，这种做法仍然合理。\par
\noindent\textbf{贡献与意义}\ 论文较早把真实排故中的认知过程描述为层级系统上的搜索，并揭示“信息最省”与“时间、心智成本最优”并不相同。\par
\clearpage
\begin{center}
{\songti\bfseries\fontsize{18.5pt}{27.75pt}\selectfont 华中科技大学设计学院硕士研究生\par}
{\songti\bfseries\fontsize{18.5pt}{27.75pt}\selectfont 文献阅读卡\par}
\end{center}
\vspace{0.25em}
\noindent\underline{\makebox[\textwidth][l]{\textbf{学号：M202570493\hfill 姓名：陈思静\hfill 序号：10}}}\par
\vspace{0.35em}
\noindent\textbf{名称}\ Making Sense of Sensemaking 1: Alternative Perspectives\par
\noindent\textbf{作者}\ Gary Klein; Brian Moon; Robert R. Hoffman\quad 译者\par
\noindent\textbf{出处}\ 【美国】IEEE Intelligent Systems，2006，21(4)：70–73\quad 页\par
\vspace{0.35em}
\noindent\textbf{文献阅读心得（每篇不少于150字）：}\par
\noindent\textbf{研究背景}\ 意义建构（sensemaking）在心理学、组织研究、自然情境决策和人机交互（human-computer interaction，HCI）中有多种含义。复杂任务中的理解不是对信息做一次摘要，而是持续寻找联系、解释意外并为下一步行动建立可用认识。\par
\noindent\textbf{研究问题}\ 研究问题一（RQ1）：不同研究传统如何界定意义建构，它与理解、情境意识等概念有何关系？\par
研究问题二（RQ2）：哪种定义能解释人在不确定环境中为了预测和行动而持续组织联系？\par
\noindent\textbf{研究方法}\ 本文是概念综述与立场论文。作者比较多个学科的定义和代表性案例，讨论意义建构的触发条件、目的与分析尺度，并为后续宏认知模型建立概念边界。\par
\noindent\textbf{主要发现}\ 作者强调意义建构是一种有动机、持续的努力：理解人、地点和事件之间的联系，以预测其发展并有效行动。它常由惊讶、冲突或知识缺口触发，结果是可行动的解释，而非静态信息表征。\par
\noindent\textbf{贡献与意义}\ 论文跨学科澄清了意义建构的不同视角，把“联系—预测—行动”作为共同核心，为研究真实复杂工作中的理解活动提供了统一入口。\par
\clearpage
\begin{center}
{\songti\bfseries\fontsize{18.5pt}{27.75pt}\selectfont 华中科技大学设计学院硕士研究生\par}
{\songti\bfseries\fontsize{18.5pt}{27.75pt}\selectfont 文献阅读卡\par}
\end{center}
\vspace{0.25em}
\noindent\underline{\makebox[\textwidth][l]{\textbf{学号：M202570493\hfill 姓名：陈思静\hfill 序号：11}}}\par
\vspace{0.35em}
\noindent\textbf{名称}\ Making Sense of Sensemaking 2: A Macrocognitive Model\par
\noindent\textbf{作者}\ Gary Klein; Brian Moon; Robert R. Hoffman\quad 译者\par
\noindent\textbf{出处}\ 【美国】IEEE Intelligent Systems，2006，21(5)：88–92\quad 页\par
\vspace{0.35em}
\noindent\textbf{文献阅读心得（每篇不少于150字）：}\par
\noindent\textbf{研究背景}\ 人在不确定任务中会借助已有框架筛选数据，数据又可能迫使框架改变。把意义建构理解成固定步骤，会忽略专家在解释、注意分配、反馈和行动之间的往返。\par
\noindent\textbf{研究问题}\ 研究问题一（RQ1）：数据与解释框架如何相互依存并推动意义建构？\par
研究问题二（RQ2）：人何时会扩展、质疑、比较、保留或重构当前框架？\par
\noindent\textbf{研究方法}\ 作者基于自然情境决策研究提出数据—框架（Data/Frame）宏认知模型，并用因果推理、假设承诺、反馈学习、专家差异和确认偏差等经验研究检验模型含义。\par
\noindent\textbf{主要发现}\ 框架管理注意并定义、连接和过滤数据；异常与不一致可触发质疑或重构。过早形成框架并非必然有害，适度承诺能产生可检验预期；过程反馈通常比只告知成败的结果反馈更有助于改进。\par
\noindent\textbf{贡献与意义}\ 论文以双向的数据—框架循环替代线性“先收集数据、后得出解释”，并提出连接、扩展、质疑、比较、重构和寻找框架等可观察活动。\par
\clearpage
\begin{center}
{\songti\bfseries\fontsize{18.5pt}{27.75pt}\selectfont 华中科技大学设计学院硕士研究生\par}
{\songti\bfseries\fontsize{18.5pt}{27.75pt}\selectfont 文献阅读卡\par}
\end{center}
\vspace{0.25em}
\noindent\underline{\makebox[\textwidth][l]{\textbf{学号：M202570493\hfill 姓名：陈思静\hfill 序号：12}}}\par
\vspace{0.35em}
\noindent\textbf{名称}\ A Data–Frame Theory of Sensemaking\par
\noindent\textbf{作者}\ Gary Klein; Jennifer K. Phillips; Erica L. Rall; Deborah A. Peluso\quad 译者\par
\noindent\textbf{出处}\ 【美国】Expertise Out of Context，Mahwah, NJ：Lawrence Erlbaum Associates，2007：113–155\quad 页\par
\vspace{0.35em}
\noindent\textbf{文献阅读心得（每篇不少于150字）：}\par
\noindent\textbf{研究背景}\ 传统微观认知研究多解释人如何理解单词、句子等简单刺激，难以覆盖真实工作中对模糊、矛盾、动态信息形成可行动理解的过程。数据—框架理论（Data–Frame theory）由此转向自然情境中的宏观意义建构。\par
\noindent\textbf{研究问题}\ 研究问题一（RQ1）：数据如何被框架选择和解释，框架又如何由少量关键锚点被识别或建构？\par
研究问题二（RQ2）：当数据与框架不一致时，意义建构有哪些可能形式与转换路径？\par
\noindent\textbf{研究方法}\ 本章综合自然情境决策、专业知识和问题解决研究，借助真实事件材料提出9条理论命题，并系统描述连接数据与框架、扩展、质疑、保留、比较、重构和寻找新框架等活动。\par
\noindent\textbf{主要发现}\ 数据不是完全中性的输入，而是在框架中被辨认和赋义；框架也由关键线索推断。意义建构通常在二者达到足够一致时停止，并服务于“接下来做什么”的功能性理解。专家与新手主要差在可调用框架的丰富度。\par
\noindent\textbf{贡献与意义}\ 理论同时说明数据驱动和框架驱动过程，明确列出多种非线性活动及其触发关系，可用于分析解释如何被建立、维护和推翻。\par
\clearpage
\begin{center}
{\songti\bfseries\fontsize{18.5pt}{27.75pt}\selectfont 华中科技大学设计学院硕士研究生\par}
{\songti\bfseries\fontsize{18.5pt}{27.75pt}\selectfont 文献阅读卡\par}
\end{center}
\vspace{0.25em}
\noindent\underline{\makebox[\textwidth][l]{\textbf{学号：M202570493\hfill 姓名：陈思静\hfill 序号：13}}}\par
\vspace{0.35em}
\noindent\textbf{名称}\ AgentGUI: An Interface for Observing and Steering Long-Running AI Agents\par
\noindent\textbf{作者}\ Xuan Zhao; Jiwoong Sohn; Qinyue Zheng; Michael Moor\quad 译者\par
\noindent\textbf{出处}\ 【瑞士】arXiv preprint arXiv:2607.26300，2026\quad 页\par
\vspace{0.35em}
\noindent\textbf{文献阅读心得（每篇不少于150字）：}\par
\noindent\textbf{研究背景}\ 长时运行的智能体可同时执行多个任务，但人类监督仍依赖终端输出和分散日志。用户需要在不中断整体工作的情况下掌握轨迹、产物、成本和偏移，并在必要时实施人工或自动纠偏。\par
\noindent\textbf{研究问题}\ 研究问题一（RQ1）：统一的轨迹可视化能否提高用户理解长时智能体执行的速度和准确性？\par
研究问题二（RQ2）：人工与自动的管理者式引导（steering）能否在任务漂移时改善最终完成率？\par
\noindent\textbf{研究方法}\ 系统整合多会话轨迹、工作区快照、人工输入和自动管理模块。8名硕博参与者在被试内、顺序平衡实验中比较AgentGUI与Hermes仪表板（Hermes Dashboard）；另对多种模型各运行50次以测试自动纠偏。\par
\noindent\textbf{主要发现}\ 参与者用AgentGUI回答轨迹问题平均快38\%（90秒对145秒，$p=.023$），准确率为93\%对80\%（$p=.031$），且没有速度—准确率交换；自动管理模块在部分模型上将完成率提高最多34个百分点。\par
\noindent\textbf{贡献与意义}\ 论文把并行会话、轨迹理解、工作区产物、人工指令和自动审计放进同一桌面式界面，并用客观理解速度、准确率和工作负荷评价可观测性。\par
\clearpage
\begin{center}
{\songti\bfseries\fontsize{18.5pt}{27.75pt}\selectfont 华中科技大学设计学院硕士研究生\par}
{\songti\bfseries\fontsize{18.5pt}{27.75pt}\selectfont 文献阅读卡\par}
\end{center}
\vspace{0.25em}
\noindent\underline{\makebox[\textwidth][l]{\textbf{学号：M202570493\hfill 姓名：陈思静\hfill 序号：14}}}\par
\vspace{0.35em}
\noindent\textbf{名称}\ An Exploratory Study of How Developers Seek, Relate, and Collect Relevant Information during Software Maintenance Tasks\par
\noindent\textbf{作者}\ Amy J. Ko; Brad A. Myers; Michael J. Coblenz; Htet Htet Aung\quad 译者\par
\noindent\textbf{出处}\ 【美国】IEEE Transactions on Software Engineering，2006，32(12)：971–987，DOI: 10.1109/TSE.2006.116\quad 页\par
\vspace{0.35em}
\noindent\textbf{文献阅读心得（每篇不少于150字）：}\par
\noindent\textbf{研究背景}\ 软件维护不仅需要读懂局部代码，还要在多文件、多工具和多任务之间寻找、关联并暂存相关信息。现代集成开发环境（integrated development environment，IDE）的导航机制可能既提供线索，也制造误导与重复切换。\par
\noindent\textbf{研究问题}\ 研究问题一（RQ1）：开发者在维护任务中如何搜索、关联和收集相关代码与运行信息？\par
研究问题二（RQ2）：集成开发环境的线索、依赖导航和交互状态如何支持或妨碍这些活动？\par
\noindent\textbf{研究方法}\ 研究从31名Java开发者中选取经验最高的10名，要求其在70分钟内完成同一503行Java Paint程序的5个维护任务，并接受定时中断。作者逐帧分析屏幕录像、语音和交互日志。\par
\noindent\textbf{主要发现}\ 开发者平均35\%的有效时间用于源文件内外的导航，许多跳转是重复且本可避免的；失败搜索常由误导性信息气味开始。记忆不可靠、纸笔缓慢、集成开发环境的交互状态易被后续任务覆盖，均难稳定保存任务上下文。\par
\noindent\textbf{贡献与意义}\ 论文提出“搜索—关联—收集”的程序理解模型，将信息觅食与集成开发环境的具体交互结合，并以真实多任务和中断条件揭示导航成本及外部记忆需求。\par
\clearpage
\begin{center}
{\songti\bfseries\fontsize{18.5pt}{27.75pt}\selectfont 华中科技大学设计学院硕士研究生\par}
{\songti\bfseries\fontsize{18.5pt}{27.75pt}\selectfont 文献阅读卡\par}
\end{center}
\vspace{0.25em}
\noindent\underline{\makebox[\textwidth][l]{\textbf{学号：M202570493\hfill 姓名：陈思静\hfill 序号：15}}}\par
\vspace{0.35em}
\noindent\textbf{名称}\ Debugging: An Analysis of Bug-Location Strategies\par
\noindent\textbf{作者}\ Irvin R. Katz; John R. Anderson\quad 译者\par
\noindent\textbf{出处}\ 【美国】Human–Computer Interaction，1987–1988，3(4)：351–399，DOI: 10.1207/S15327051HCI0304\_2\quad 页\par
\vspace{0.35em}
\noindent\textbf{文献阅读心得（每篇不少于150字）：}\par
\noindent\textbf{研究背景}\ 具备语言知识的程序员往往能在看到错误行后修正它，真正困难的是从程序行为中定位哪一行有错。调试策略可能因代码归属、可用线索和搜索顺序而改变。\par
\noindent\textbf{研究问题}\ 研究问题一（RQ1）：熟悉表处理语言（LISP）的学生在定位程序错误时采用哪些策略，错误更接近知识缺陷还是偶发失误？\par
研究问题二（RQ2）：调试自己的程序与他人程序是否改变定位策略，策略又如何影响检查与修复？\par
\noindent\textbf{研究方法}\ 论文报告4项表处理语言调试实验，比较错误重复性、自然使用的定位策略、自编与他编程序情境，以及不同策略对被检查代码行和后续判断修复的影响。\par
\noindent\textbf{主要发现}\ 参与者的错误多为不易重复的slip，通常一旦定位即可修复；他们会使用多种错误定位策略，且自编与他编程序的策略选择不同。策略主要改变搜索哪些代码行，对选中一行后的正确性判断和修复能力影响较小。\par
\noindent\textbf{贡献与意义}\ 研究把“找错位置”与“判断并修复错误”分解为不同阶段，并通过连续实验说明调试瓶颈并非总是缺少编程知识，而是搜索与定位。\par
\clearpage
\begin{center}
{\songti\bfseries\fontsize{18.5pt}{27.75pt}\selectfont 华中科技大学设计学院硕士研究生\par}
{\songti\bfseries\fontsize{18.5pt}{27.75pt}\selectfont 文献阅读卡\par}
\end{center}
\vspace{0.25em}
\noindent\underline{\makebox[\textwidth][l]{\textbf{学号：M202570493\hfill 姓名：陈思静\hfill 序号：16}}}\par
\vspace{0.35em}
\noindent\textbf{名称}\ Theory of Troubleshooting: The Developer’s Cognitive Experience of Overcoming Confusion\par
\noindent\textbf{作者}\ Arty Starr; Margaret-Anne Storey\quad 译者\par
\noindent\textbf{出处}\ 【美国】ACM Transactions on Software Engineering and Methodology，2026，DOI: 10.1145/3800945；arXiv:2602.10540\quad 页\par
\vspace{0.35em}
\noindent\textbf{文献阅读心得（每篇不少于150字）：}\par
\noindent\textbf{研究背景}\ 调试研究常聚焦工具和步骤，却较少解释开发者面对意外行为时的内部体验。长时间困惑会持续占用注意、工作记忆和心理模型构建资源，并可能转化为疲劳与项目可持续性风险。\par
\noindent\textbf{研究问题}\ 研究问题一（RQ1）：开发者在排查意外系统行为时，思考、感受和努力随时间如何变化？\par
研究问题二（RQ2）：困惑、经验直觉、试探性实验、进展与认知疲劳之间存在怎样的关系？\par
\noindent\textbf{研究方法}\ 作者访谈27名职业开发者，合计570年经验，并开展7次回访；采用建构主义扎根理论，从681条扎根编码形成8个核心理论类别，在前12份访谈后达到理论饱和，并以参与者反馈检验共鸣。\par
\noindent\textbf{主要发现}\ 排故从困惑体验启动，经验直觉引导策略，开发者通过“试一试、看结果”缩小解释空间；获得正或负知识都可被体验为进展。持续无进展会耗竭注意并引发补偿努力、盲视和疲劳，最终在“想通了”时解除。\par
\noindent\textbf{贡献与意义}\ 论文把排故定义为识别、理解并建立意外行为原因心理模型的认知过程，用连贯理论连接情绪、注意、实验、直觉和疲劳，而不只描述技术操作。\par
\clearpage
\begin{center}
{\songti\bfseries\fontsize{18.5pt}{27.75pt}\selectfont 华中科技大学设计学院硕士研究生\par}
{\songti\bfseries\fontsize{18.5pt}{27.75pt}\selectfont 文献阅读卡\par}
\end{center}
\vspace{0.25em}
\noindent\underline{\makebox[\textwidth][l]{\textbf{学号：M202570493\hfill 姓名：陈思静\hfill 序号：17}}}\par
\vspace{0.35em}
\noindent\textbf{名称}\ The Sensemaking Process and Leverage Points for Analyst Technology as Identified Through Cognitive Task Analysis\par
\noindent\textbf{作者}\ Peter Pirolli; Stuart K. Card\quad 译者\par
\noindent\textbf{出处}\ 【美国】Proceedings of the 2005 International Conference on Intelligence Analysis，2005\quad 页\par
\vspace{0.35em}
\noindent\textbf{文献阅读心得（每篇不少于150字）：}\par
\noindent\textbf{研究背景}\ 情报分析需要从大量异质材料形成结构、假设和可交付结论，但规范性流程很少说明专家真实工作中如何往返搜索、筛选、重组和验证证据。\par
\noindent\textbf{研究问题}\ 研究问题一（RQ1）：专家分析者如何在信息觅食与意义建构之间循环，将信息转化为模式、洞见和产品？\par
研究问题二（RQ2）：认知任务分析揭示了哪些可由交互技术降低成本的关键杠杆点？\par
\noindent\textbf{研究方法}\ 作者采用认知任务分析、分析者访谈和出声思考协议，结合专家图式与信息觅食理论，构建从外部数据源到证据文件、图式、假设和最终报告的双环模型。\par
\noindent\textbf{主要发现}\ 分析并非线性阅读，而是在探索、富集、利用与表示重构之间机会主义往返；常见外部结构包括资料箱（shoebox）、证据文件（evidence file）和图式（schema）。高成本环节集中在搜索、筛选、重新表示、比较假设和组装证据。\par
\noindent\textbf{贡献与意义}\ 论文将信息觅食循环（foraging loop）与意义建构循环（sensemaking loop）连接起来，并从每个转换的时间和注意成本推导工具设计杠杆，为总览—筛选—证据—假设界面提供理论基础。\par
\clearpage
\begin{center}
{\songti\bfseries\fontsize{18.5pt}{27.75pt}\selectfont 华中科技大学设计学院硕士研究生\par}
{\songti\bfseries\fontsize{18.5pt}{27.75pt}\selectfont 文献阅读卡\par}
\end{center}
\vspace{0.25em}
\noindent\underline{\makebox[\textwidth][l]{\textbf{学号：M202570493\hfill 姓名：陈思静\hfill 序号：18}}}\par
\vspace{0.35em}
\noindent\textbf{名称}\ AgentStepper: Interactive Debugging of Software Development Agents\par
\noindent\textbf{作者}\ Robert Hutter; Michael Pradel\quad 译者\par
\noindent\textbf{出处}\ 【德国】arXiv preprint arXiv:2602.06593，2026\quad 页\par
\vspace{0.35em}
\noindent\textbf{文献阅读心得（每篇不少于150字）：}\par
\noindent\textbf{研究背景}\ 软件开发智能体会在模型请求、工具调用和代码修改之间形成长轨迹。原始日志虽然保留了大量信息，却缺少适合开发者理解高层行为的结构，也不能像传统调试器那样暂停、单步和检查中间状态。\par
\noindent\textbf{研究问题}\ 研究问题一（RQ1）：将传统调试中的断点、单步执行和状态检查提升到智能体动作层，能否帮助开发者理解智能体轨迹？\par
研究问题二（RQ2）：AgentStepper接入现有软件开发智能体需要多少改动，它是否改善轨迹理解、智能体程序缺陷识别和主观负荷？\par
\noindent\textbf{研究方法}\ 作者把AgentStepper接入ExecutionAgent、SWE-Agent和RepairAgent，统计集成改动量；随后邀请12名参与者完成一项轨迹理解任务和两项智能体调试任务，比较AgentStepper与自行选择的常规日志工具。\par
\noindent\textbf{主要发现}\ 三种智能体各需插入5—7个应用程序编程接口（application programming interface，API）调用并修改39—42行代码。AgentStepper条件下轨迹理解平均表现由64\%升至67\%，智能体程序缺陷识别成功率由17\%升至60\%，挫败感由5.4/7降至2.4/7。\par
\noindent\textbf{贡献与意义}\ 论文把智能体轨迹表示为“模型—智能体程序”和“智能体程序—工具”两条交织会话，并整合总览、细节、代码变更和断点控制。它是本研究原型最直接的人侧调试先例，也说明“理解轨迹”和“定位运行框架实现缺陷”应分开测量。\par
\clearpage
\begin{center}
{\songti\bfseries\fontsize{18.5pt}{27.75pt}\selectfont 华中科技大学设计学院硕士研究生\par}
{\songti\bfseries\fontsize{18.5pt}{27.75pt}\selectfont 文献阅读卡\par}
\end{center}
\vspace{0.25em}
\noindent\underline{\makebox[\textwidth][l]{\textbf{学号：M202570493\hfill 姓名：陈思静\hfill 序号：19}}}\par
\vspace{0.35em}
\noindent\textbf{名称}\ AgentSight: System-Level Observability for AI Agents Using eBPF\par
\noindent\textbf{作者}\ Yusheng Zheng; Yanpeng Hu; Tong Yu; Andi Quinn\quad 译者\par
\noindent\textbf{出处}\ 【美国】arXiv preprint arXiv:2508.02736，2025\quad 页\par
\vspace{0.35em}
\noindent\textbf{文献阅读心得（每篇不少于150字）：}\par
\noindent\textbf{研究背景}\ 智能体的高层意图存在于模型提示和响应中，实际效果却表现为进程、文件和网络事件。现有工具通常只能看到其中一侧，导致开发者无法判断某个系统行为是否由智能体意图、工具实现或外部环境造成。\par
\noindent\textbf{研究问题}\ 研究问题一（RQ1）：能否在不修改智能体框架代码的情况下关联模型层语义意图与系统层实际行为？\par
研究问题二（RQ2）：这种跨层观测能否以可接受的开销识别安全攻击、低效循环和多智能体协调问题？\par
\noindent\textbf{研究方法}\ 作者提出边界追踪（boundary tracing），利用扩展伯克利包过滤器（extended Berkeley Packet Filter，eBPF）监视内核事件，同时截获传输层安全协议（Transport Layer Security，TLS）加密的模型流量以提取语义意图，再通过实时关联引擎和辅助大语言模型跨进程连接两类事件，并在多种攻击和故障场景中评估。\par
\noindent\textbf{主要发现}\ AgentSight能够检测提示注入、识别浪费资源的推理循环并发现隐藏的多智能体协调瓶颈，运行开销低于3\%。研究也表明，系统级事实和语义解释来自不同观测通道。\par
\noindent\textbf{贡献与意义}\ 论文具体展示了跨层证据如何被关联，并揭示“直接记录的系统事件”和“经过模型分析得到的语义关系”不能混为同一种证据。它为本研究原型的记录、推导和未知三类证据状态提供了技术边界。\par
\clearpage
\begin{center}
{\songti\bfseries\fontsize{18.5pt}{27.75pt}\selectfont 华中科技大学设计学院硕士研究生\par}
{\songti\bfseries\fontsize{18.5pt}{27.75pt}\selectfont 文献阅读卡\par}
\end{center}
\vspace{0.25em}
\noindent\underline{\makebox[\textwidth][l]{\textbf{学号：M202570493\hfill 姓名：陈思静\hfill 序号：20}}}\par
\vspace{0.35em}
\noindent\textbf{名称}\ From Failed Trajectories to Reliable LLM Agents: Diagnosing and Repairing Harness Flaws\par
\noindent\textbf{作者}\ Mengzhuo Chen; Junjie Wang; Zhe Liu; Yawen Wang; Haiming Zheng; Qing Wang\quad 译者\par
\noindent\textbf{出处}\ 【中国】arXiv preprint arXiv:2606.06324v2，2026\quad 页\par
\vspace{0.35em}
\noindent\textbf{文献阅读心得（每篇不少于150字）：}\par
\noindent\textbf{研究背景}\ 智能体失败往往由执行环境、工具、上下文、生命周期、可观测、验证或治理机制共同塑造。仅依据最终成败调整提示词，难以确定失败证据位于哪一步，也难以把运行现象映射到具体运行框架工件。\par
\noindent\textbf{研究问题}\ 研究问题一（RQ1）：如何把碎片化运行轨迹编译成能够表达步骤级数据流、控制流和运行框架实现锚点的统一表示？\par
研究问题二（RQ2）：基于该表示形成的局部诊断和范围化修复，能否在不引入不可接受回归的情况下改善智能体表现？\par
\noindent\textbf{研究方法}\ HarnessFix将原始轨迹与运行框架工件编译为面向运行框架的轨迹中间表示（Harness-aware Trace Intermediate Representation，HTIR），定位责任步骤和实现工件，把重复诊断合并为缺陷记录，再按缺陷类型生成补丁并做回归约束验证；实验覆盖四个公开基准。\par
\noindent\textbf{主要发现}\ 相较初始运行框架，HarnessFix在四个基准上的性能提高6.3至18.4个百分点，并优于人工设计和自演化基线。结果显示，失败轨迹可被组织成指向具体运行框架机制的结构化证据。\par
\noindent\textbf{贡献与意义}\ 论文证明数据流、控制流和实现锚点可以同时进入面向运行框架的轨迹表示，是与本研究原型最接近的技术工作。其自动修复不是本研究的目标，但HTIR所物化的关联类型构成了人侧理解界面的关键比较对象。\par
\clearpage
\begin{center}
{\songti\bfseries\fontsize{18.5pt}{27.75pt}\selectfont 华中科技大学设计学院硕士研究生\par}
{\songti\bfseries\fontsize{18.5pt}{27.75pt}\selectfont 文献阅读卡\par}
\end{center}
\vspace{0.25em}
\noindent\underline{\makebox[\textwidth][l]{\textbf{学号：M202570493\hfill 姓名：陈思静\hfill 序号：21}}}\par
\vspace{0.35em}
\noindent\textbf{名称}\ Agentic Harness Engineering: Observability-Driven Automatic Evolution of Coding-Agent Harnesses\par
\noindent\textbf{作者}\ Jiahang Lin; Shichun Liu; Chengjun Pan; Lizhi Lin; Shihan Dou; Zhiheng Xi; Xuanjing Huang; Hang Yan; Zhenhua Han; Tao Gui; Yu-Gang Jiang\quad 译者\par
\noindent\textbf{出处}\ 【中国】arXiv preprint arXiv:2604.25850v4，2026\quad 页\par
\vspace{0.35em}
\noindent\textbf{文献阅读心得（每篇不少于150字）：}\par
\noindent\textbf{研究背景}\ 智能体运行框架包含工具、中间件、长期记忆、系统提示和运行规则，修改空间异构，轨迹规模巨大，而且一次修改与下一轮结果之间的关系难以归因。自动演化若缺少可观测证据，容易退化为试错。\par
\noindent\textbf{研究问题}\ 研究问题一（RQ1）：组件、运行经验和修改决策需要怎样的可观测表示，才能形成可验证的运行框架演化闭环？\par
研究问题二（RQ2）：可观测性驱动的自动演化能否提高编码智能体性能，并跨模型迁移而非只适配单一基准？\par
\noindent\textbf{研究方法}\ 作者提出组件可观测、经验可观测和决策可观测三支柱，在Terminal-Bench 2上执行10轮运行框架演化，并通过消融和跨模型、跨基准迁移检验各类组件修改的作用。\par
\noindent\textbf{主要发现}\ Terminal-Bench 2的单次通过率（pass@1）由69.7\%升至77.0\%；冻结后的运行框架在其他模型家族上带来5.1—10.1个百分点提升，并以比种子运行框架少12\%的词元（token）取得SWE-bench Verified上的领先结果。增益主要来自工具、中间件和长期记忆，而非系统提示。\par
\noindent\textbf{贡献与意义}\ 论文把可观测性从日志采集提升为组件、经验和决策三个层次，并要求每次修改都有可证伪预测。它说明本研究不能只研究提示词，而应将上下文和工具等机制放入统一运行框架结构。\par
\clearpage
\begin{center}
{\songti\bfseries\fontsize{18.5pt}{27.75pt}\selectfont 华中科技大学设计学院硕士研究生\par}
{\songti\bfseries\fontsize{18.5pt}{27.75pt}\selectfont 文献阅读卡\par}
\end{center}
\vspace{0.25em}
\noindent\underline{\makebox[\textwidth][l]{\textbf{学号：M202570493\hfill 姓名：陈思静\hfill 序号：22}}}\par
\vspace{0.35em}
\noindent\textbf{名称}\ AgentLens: Visual Analysis for Agent Behaviors in LLM-Based Autonomous Systems\par
\noindent\textbf{作者}\ Jiaying Lu; Bo Pan; Jieyi Chen; Yingchaojie Feng; Jingyuan Hu; Yuchen Peng; Wei Chen\quad 译者\par
\noindent\textbf{出处}\ 【美国】IEEE Transactions on Visualization and Computer Graphics，2025，31(8)：4182–4197，DOI: 10.1109/TVCG.2024.3394053\quad 页\par
\vspace{0.35em}
\noindent\textbf{文献阅读心得（每篇不少于150字）：}\par
\noindent\textbf{研究背景}\ 多智能体自治系统产生大量随时间演化的事件、记忆和交互。线性日志难以呈现跨时间行为模式，也难以让用户从总体演化下钻到单个智能体行为及其可能原因。\par
\noindent\textbf{研究问题}\ 研究问题一（RQ1）：如何把原始执行事件组织成支持多粒度查看的智能体行为模型，并追踪行为之间的潜在因果关系？\par
研究问题二（RQ2）：层级时间视图、单智能体细节视图和环境监视视图能否支持用户发现并解释复杂行为？\par
\noindent\textbf{研究方法}\ 作者建立事件处理、行为摘要和原因追踪流程，开发包含概览（Outline）、智能体（Agent）和监视（Monitor）三类联动视图的AgentLens，并通过两个使用场景和用户研究评价系统。\par
\noindent\textbf{主要发现}\ 参与者能够沿层级时间线发现行为模式、下钻检查事件，并回看可能影响当前行为的历史操作。系统有效性和可用性得到验证，但部分因果关系依据文本相似度推断，并非直接观测事实。\par
\noindent\textbf{贡献与意义}\ 论文提供总览—细节联动和原因追踪的成熟视觉设计，也暴露了本研究原型要解决的风险：界面若不说明连线来自日志、算法推断还是缺失信号，使用者可能把潜在关系理解为实际因果。\par
\clearpage
\begin{center}
{\songti\bfseries\fontsize{18.5pt}{27.75pt}\selectfont 华中科技大学设计学院硕士研究生\par}
{\songti\bfseries\fontsize{18.5pt}{27.75pt}\selectfont 文献阅读卡\par}
\end{center}
\vspace{0.25em}
\noindent\underline{\makebox[\textwidth][l]{\textbf{学号：M202570493\hfill 姓名：陈思静\hfill 序号：23}}}\par
\vspace{0.35em}
\noindent\textbf{名称}\ AgentGraph: Trace-to-Graph Platform for Interactive Analysis and Robustness Testing in Agentic AI Systems\par
\noindent\textbf{作者}\ Zekun Wu; Seonglae Cho; Cristian Enrique Munoz Villalobos; Theo King; Umar Mohammed; Emre Kazim; Maria Perez-Ortiz; Sahan Bulathwela; Adriano Koshiyama\quad 译者\par
\noindent\textbf{出处}\ 【美国】Proceedings of the AAAI Conference on Artificial Intelligence，2026，40(48)：41721–41723，DOI: 10.1609/aaai.v40i48.42393\quad 页\par
\vspace{0.35em}
\noindent\textbf{文献阅读心得（每篇不少于150字）：}\par
\noindent\textbf{研究背景}\ 现有智能体可观测平台多记录提示输入输出和运行指标，却仍要求用户手工从线性轨迹恢复任务、工具、数据和人员之间的结构关系。\par
\noindent\textbf{研究问题}\ 研究问题一（RQ1）：如何把智能体执行日志转换为带类型节点和边的交互知识图，并保持每个图元素可回到原始轨迹？\par
研究问题二（RQ2）：统一图表示能否同时支持故障分析、优化建议和扰动条件下的稳健性测试？\par
\noindent\textbf{研究方法}\ AgentGraph把智能体、任务、工具、数据输入输出和人建模为节点，用消费、委托、顺序、工具使用、输出交付和人工干预等类型边连接，并让每个元素指向精确轨迹片段（trace span）；系统还加入扰动和因果归因分析。\par
\noindent\textbf{主要发现}\ 原型能够在同一图中开展基于轨迹的定性故障分析和定量稳健性测试。论文属于演示轨道，主要证明系统能力，尚未提供充分的人类理解效果证据。\par
\noindent\textbf{贡献与意义}\ 论文将关系图与精确原始证据链接，是本研究原型“关联必须可核验”的直接先例；同时也说明系统功能展示不能替代对理解正确率、误判率和无证据结论的受控评估。\par
\clearpage
\begin{center}
{\songti\bfseries\fontsize{18.5pt}{27.75pt}\selectfont 华中科技大学设计学院硕士研究生\par}
{\songti\bfseries\fontsize{18.5pt}{27.75pt}\selectfont 文献阅读卡\par}
\end{center}
\vspace{0.25em}
\noindent\underline{\makebox[\textwidth][l]{\textbf{学号：M202570493\hfill 姓名：陈思静\hfill 序号：24}}}\par
\vspace{0.35em}
\noindent\textbf{名称}\ Graph of Trace: Visualizing Execution Traces of Scientific Agents\par
\noindent\textbf{作者}\ Tianci Gao; Haoxuan Li; Jian He Li; Tianxiang Zhao; Shi Runze; Weiran Wang; Zezhao Wu; Lu Mi\quad 译者\par
\noindent\textbf{出处}\ 【美国】Proceedings of the 64th Annual Meeting of the Association for Computational Linguistics: System Demonstrations，2026：297–306，DOI: 10.18653/v1/2026.acl-demo.29\quad 页\par
\vspace{0.35em}
\noindent\textbf{文献阅读心得（每篇不少于150字）：}\par
\noindent\textbf{研究背景}\ 科学智能体可以自动执行环境配置、代码运行、数据分析和结论生成，但展开后的工作流难以供领域专家检查，限制了结果复核和后续人机协作。\par
\noindent\textbf{研究问题}\ 研究问题一（RQ1）：如何实时记录细粒度执行事件并把科学智能体工作流组织成可检查的有向图？\par
研究问题二（RQ2）：结构化轨迹图是否改善领域专家对工作流、失败位置和结果生成过程的理解？\par
\noindent\textbf{研究方法}\ 系统记录工具调用、代码执行和中间产物，由监视模块构建实时更新的有向图，并提供对话、图和节点细节检查器；作者在复杂科研任务上邀请人工智能、神经科学和生物学专家评价。\par
\noindent\textbf{主要发现}\ 专家认为结构化轨迹提高了工作流理解、感知可解释性和进一步分析交互的可用性。评价以专家反馈为主，尚未通过同证据基线比较理解准确率和无依据判断。\par
\noindent\textbf{贡献与意义}\ 论文提供与本研究原型相近的实时轨迹图和总览—节点细节结构，同时构成重要方法学对照：本研究需要把主观可用性扩展为具有客观真值（ground truth）的理解和证据判断实验。\par
\clearpage
\begin{center}
{\songti\bfseries\fontsize{18.5pt}{27.75pt}\selectfont 华中科技大学设计学院硕士研究生\par}
{\songti\bfseries\fontsize{18.5pt}{27.75pt}\selectfont 文献阅读卡\par}
\end{center}
\vspace{0.25em}
\noindent\underline{\makebox[\textwidth][l]{\textbf{学号：M202570493\hfill 姓名：陈思静\hfill 序号：25}}}\par
\vspace{0.35em}
\noindent\textbf{名称}\ AUTOGEN STUDIO: A No-Code Developer Tool for Building and Debugging Multi-Agent Systems\par
\noindent\textbf{作者}\ Victor Dibia; Jingya Chen; Gagan Bansal; Suff Syed; Adam Fourney; Erkang Zhu; Chi Wang; Saleema Amershi\quad 译者\par
\noindent\textbf{出处}\ 【美国】Proceedings of the 2024 Conference on Empirical Methods in Natural Language Processing: System Demonstrations，2024：72–79，DOI: 10.18653/v1/2024.emnlp-demo.8\quad 页\par
\vspace{0.35em}
\noindent\textbf{文献阅读心得（每篇不少于150字）：}\par
\noindent\textbf{研究背景}\ 多智能体工作流涉及模型、工具、角色和编排参数，直接通过代码配置和日志排查门槛较高。开发者需要同时查看系统声明结构和具体运行结果。\par
\noindent\textbf{研究问题}\ 研究问题一（RQ1）：声明式规范和无代码界面如何降低多智能体工作流构建、调试和评价的复杂度？\par
研究问题二（RQ2）：哪些界面原则能够让模型、工具、智能体和编排关系被复用并在运行中检查？\par
\noindent\textbf{研究方法}\ 作者基于AutoGen开发网页（Web）界面和Python应用程序编程接口，以JavaScript对象表示法（JavaScript Object Notation，JSON）声明智能体及工作流，提供拖放式结构编辑、交互运行与调试、结果评估和可复用组件库，并通过系统案例说明设计。\par
\noindent\textbf{主要发现}\ AutoGen Studio把工作流结构、组件配置和运行查看统一到一个开发环境中，支持快速原型和非代码式调试。论文主要是系统演示，没有隔离不同呈现方式对理解正确率的影响。\par
\noindent\textbf{贡献与意义}\ 论文证明“声明结构”和“运行轨迹”可以在同一开发工具中组织，并提供可复用的智能体系统表达方式；本研究进一步关心静态声明与本次运行事实之间的证据差异。\par
\clearpage
\begin{center}
{\songti\bfseries\fontsize{18.5pt}{27.75pt}\selectfont 华中科技大学设计学院硕士研究生\par}
{\songti\bfseries\fontsize{18.5pt}{27.75pt}\selectfont 文献阅读卡\par}
\end{center}
\vspace{0.25em}
\noindent\underline{\makebox[\textwidth][l]{\textbf{学号：M202570493\hfill 姓名：陈思静\hfill 序号：26}}}\par
\vspace{0.35em}
\noindent\textbf{名称}\ AgentTrace: A Structured Logging Framework for Agent System Observability\par
\noindent\textbf{作者}\ Adam AlSayyad; Kelvin Yuxiang Huang; Richik Pal\quad 译者\par
\noindent\textbf{出处}\ 【美国】arXiv preprint arXiv:2602.10133，2026\quad 页\par
\vspace{0.35em}
\noindent\textbf{文献阅读心得（每篇不少于150字）：}\par
\noindent\textbf{研究背景}\ 传统日志偏向操作事件，难以同时表达智能体的模型交互、环境动作和上下文状态。静态审计也难以应对由模型推断产生的非确定行为。\par
\noindent\textbf{研究问题}\ 研究问题一（RQ1）：智能体运行时日志应覆盖哪些观测面，才能支持调试、安全审计和责任追踪？\par
研究问题二（RQ2）：结构化日志能否在较少侵入智能体代码的情况下接入现有遥测基础设施？\par
\noindent\textbf{研究方法}\ 作者提出操作（operational）、认知（cognitive）和上下文（contextual）三类日志面，设计图式化（schema-based）运行时记录机制，并展示其与开放遥测框架（OpenTelemetry）等基础设施的集成方式。\par
\noindent\textbf{主要发现}\ 框架能够持续记录工具和环境动作、模型相关状态以及上下文信息，为实时监控和事后检查提供共同结构。论文以框架设计为主，尚未验证开发者如何理解这些日志。\par
\noindent\textbf{贡献与意义}\ 论文把上下文状态明确列为与操作和认知并列的观测面，直接支持本研究聚焦上下文构建与变换；同时提醒研究区分“记录结构完整”与“人能正确理解”两类贡献。\par
\clearpage
\begin{center}
{\songti\bfseries\fontsize{18.5pt}{27.75pt}\selectfont 华中科技大学设计学院硕士研究生\par}
{\songti\bfseries\fontsize{18.5pt}{27.75pt}\selectfont 文献阅读卡\par}
\end{center}
\vspace{0.25em}
\noindent\underline{\makebox[\textwidth][l]{\textbf{学号：M202570493\hfill 姓名：陈思静\hfill 序号：27}}}\par
\vspace{0.35em}
\noindent\textbf{名称}\ From Agent Traces to Trust: Evidence Tracing and Execution Provenance in LLM Agents\par
\noindent\textbf{作者}\ Yiqi Wang; Jiaqi Zhang; Taotao Cai; Zirui Liu; Qingqiang Sun; Zequn Sun; Zhangkai Wu; Mingkai Zhang; Yanming Zhu\quad 译者\par
\noindent\textbf{出处}\ 【中国】arXiv preprint arXiv:2606.04990，2026\quad 页\par
\vspace{0.35em}
\noindent\textbf{文献阅读心得（每篇不少于150字）：}\par
\noindent\textbf{研究背景}\ 最终答案正确与否不能说明证据如何进入结论、工具调用是否合理、记忆怎样影响后续决策，也不能定位执行失败。相关工作分散在检索归因、工具安全、记忆、可观测和调试领域。\par
\noindent\textbf{研究问题}\ 研究问题一（RQ1）：大语言模型智能体中的证据单元、执行单元和溯源关系可以怎样统一分类？\par
研究问题二（RQ2）：现有方法在表示、粒度、时间、评价和信任功能上有哪些覆盖与缺口？\par
\noindent\textbf{研究方法}\ 作者开展系统性综述，从轨迹（trace）来源、证据与执行单元、溯源（provenance）关系、追踪粒度与时机、表示形式和信任功能建立统一分类，并映射数据集、基准和评价指标。\par
\noindent\textbf{主要发现}\ 现有研究能够追踪检索、工具、记忆和动作的部分来源，但缺少统一的轨迹图式（trace schema）、语义级溯源、真实执行基准和面向恢复的过程评价，隐私与审计基础设施也不成熟。\par
\noindent\textbf{贡献与意义}\ 论文直接把证据追踪与执行溯源连接起来，为本研究定义“关联证据”提供最新综述框架，也支持把评价从最终答案扩展到过程可核验性和证据边界。\par
\clearpage
\begin{center}
{\songti\bfseries\fontsize{18.5pt}{27.75pt}\selectfont 华中科技大学设计学院硕士研究生\par}
{\songti\bfseries\fontsize{18.5pt}{27.75pt}\selectfont 文献阅读卡\par}
\end{center}
\vspace{0.25em}
\noindent\underline{\makebox[\textwidth][l]{\textbf{学号：M202570493\hfill 姓名：陈思静\hfill 序号：28}}}\par
\vspace{0.35em}
\noindent\textbf{名称}\ LADYBUG: an LLM Agent DeBUGger for Data-Driven Applications\par
\noindent\textbf{作者}\ Joel Rorseth; Parke Godfrey; Lukasz Golab; Divesh Srivastava; Jarek Szlichta\quad 译者\par
\noindent\textbf{出处}\ 【荷兰】Proceedings of the 28th International Conference on Extending Database Technology，2025：1082–1085，DOI: 10.48786/EDBT.2025.94\quad 页\par
\vspace{0.35em}
\noindent\textbf{文献阅读心得（每篇不少于150字）：}\par
\noindent\textbf{研究背景}\ 数据驱动智能体常以多步工作流生成代码、转换数据并汇总结果。某一步的隐蔽错误可能不触发异常，却污染后续结果，用户需要检查中间输入输出并限定重新执行范围。\par
\noindent\textbf{研究问题}\ 研究问题一（RQ1）：交互式时间线能否支持用户检查智能体每一步的输入、输出和代码，并在任意步骤介入？\par
研究问题二（RQ2）：大语言模型辅助调试器如何在用户难以发现隐蔽错误时提出候选错误步骤和干预？\par
\noindent\textbf{研究方法}\ 作者实现LADYBUG演示系统，把智能体执行表示为按时间排列的步骤，允许用户查看中间状态、修改任意步骤并只重算受影响部分，同时加入基于自我反思的大语言模型辅助调试器。\par
\noindent\textbf{主要发现}\ 系统能在示例数据应用中定位并修正未直接报错的中间错误，并通过依赖关系减少重算范围。论文是演示短文，没有人类对照实验，也不能把大语言模型提出的原因视为已证实事实。\par
\noindent\textbf{贡献与意义}\ 论文展示了步骤、依赖和局部重算如何共同支持智能体调试。对本研究而言，其更重要的意义是区分可直接检查的步骤事实、系统计算的依赖关系和大语言模型提出的候选解释。\par
\clearpage
\begin{center}
{\songti\bfseries\fontsize{18.5pt}{27.75pt}\selectfont 华中科技大学设计学院硕士研究生\par}
{\songti\bfseries\fontsize{18.5pt}{27.75pt}\selectfont 文献阅读卡\par}
\end{center}
\vspace{0.25em}
\noindent\underline{\makebox[\textwidth][l]{\textbf{学号：M202570493\hfill 姓名：陈思静\hfill 序号：29}}}\par
\vspace{0.35em}
\noindent\textbf{名称}\ Open Agent Specification: Enabling Cross-Framework Comparison of AI Agents\par
\noindent\textbf{作者}\ Soufiane Amini; Yassine Benajiba; Cesare Bernardis; Paul Cayet; et al.\quad 译者\par
\noindent\textbf{出处}\ 【美国】Proceedings of the ACM Conference on AI and Agentic Systems，2026：404–418，DOI: 10.1145/3786335.3813130\quad 页\par
\vspace{0.35em}
\noindent\textbf{文献阅读心得（每篇不少于150字）：}\par
\noindent\textbf{研究背景}\ 不同智能体框架具有不同的提示处理、工具适配、编排和执行语义。即使智能体设计看似相同，运行时也可能产生不同路径和结果，导致跨框架比较缺少共同参照。\par
\noindent\textbf{研究问题}\ 研究问题一（RQ1）：如何用统一声明表示固定的智能体结构、数据流和控制流，并在不同运行时中生成可比执行？\par
研究问题二（RQ2）：相同智能体规范在不同框架中是否仍会出现准确率、时延和执行行为差异？\par
\noindent\textbf{研究方法}\ 作者提出框架无关的智能体规范（Agent Specification）及标准化评估运行框架，在LangGraph、CrewAI、AutoGen和WayFlow四种运行时上执行SimpleQA Verified、$\tau^2$-Bench和BIRD-SQL三个基准，并记录统一结构的运行轨迹（trace）。\par
\noindent\textbf{主要发现}\ 在共享规范下，不同运行时仍表现出显著的准确率、时延和执行差异，说明框架适配器、提示和编排实现会塑造运行行为，不能由静态规范直接推断。\par
\noindent\textbf{贡献与意义}\ 论文为结构与运行的对照提供了统一表示和跨框架证据，直接支撑本研究区分“预设结构”与“本次实际执行”；它也为构造相同任务、不同运行框架机制的实验材料提供方法。\par
\clearpage
\begin{center}
{\songti\bfseries\fontsize{18.5pt}{27.75pt}\selectfont 华中科技大学设计学院硕士研究生\par}
{\songti\bfseries\fontsize{18.5pt}{27.75pt}\selectfont 文献阅读卡\par}
\end{center}
\vspace{0.25em}
\noindent\underline{\makebox[\textwidth][l]{\textbf{学号：M202570493\hfill 姓名：陈思静\hfill 序号：30}}}\par
\vspace{0.35em}
\noindent\textbf{名称}\ X-Trace: A Pervasive Network Tracing Framework\par
\noindent\textbf{作者}\ Rodrigo Fonseca; George Porter; Randy H. Katz; Scott Shenker; Ion Stoica\quad 译者\par
\noindent\textbf{出处}\ 【美国】Proceedings of the 4th USENIX Symposium on Networked Systems Design and Implementation，2007：271–284\quad 页\par
\vspace{0.35em}
\noindent\textbf{文献阅读心得（每篇不少于150字）：}\par
\noindent\textbf{研究背景}\ 复杂网络服务跨越应用、协议和管理域，单层日志不能重建一次用户任务经过的完整路径。即使各组件都产生日志，也缺少将分散事件关联到同一任务的机制。\par
\noindent\textbf{研究问题}\ 研究问题一（RQ1）：如何在跨层、跨应用执行中传播统一任务标识并重建因果任务树？\par
研究问题二（RQ2）：该机制在DNS、Web和覆盖网络等异构系统中能否提供可用于诊断的端到端视图？\par
\noindent\textbf{研究方法}\ X-Trace在初始任务中注入元数据，并沿协议调用、转发和递归请求传播；各组件记录带共同任务标识的事件，系统据此重建任务树，并在DNS、三层网站和覆盖网络场景中部署。\par
\noindent\textbf{主要发现}\ 跨组件任务树能把原本分散的日志连接起来并定位错误数据或异常路径；即使只在部分层部署，仍可提供局部诊断价值，但未被记录的分支不能被当作未发生。\par
\noindent\textbf{贡献与意义}\ 论文奠定了端到端执行溯源的任务树思想，说明可靠关联依赖显式传播的记录机制。它为本研究界定“直接记录”提供技术基础，也说明信号缺失应标为未知。\par
\clearpage
\begin{center}
{\songti\bfseries\fontsize{18.5pt}{27.75pt}\selectfont 华中科技大学设计学院硕士研究生\par}
{\songti\bfseries\fontsize{18.5pt}{27.75pt}\selectfont 文献阅读卡\par}
\end{center}
\vspace{0.25em}
\noindent\underline{\makebox[\textwidth][l]{\textbf{学号：M202570493\hfill 姓名：陈思静\hfill 序号：31}}}\par
\vspace{0.35em}
\noindent\textbf{名称}\ Pivot Tracing: Dynamic Causal Monitoring for Distributed Systems\par
\noindent\textbf{作者}\ Jonathan Mace; Ryan Roelke; Rodrigo Fonseca\quad 译者\par
\noindent\textbf{出处}\ 【美国】Proceedings of the 25th ACM Symposium on Operating Systems Principles，2015：378–393，DOI: 10.1145/2815400.2815415\quad 页\par
\vspace{0.35em}
\noindent\textbf{文献阅读心得（每篇不少于150字）：}\par
\noindent\textbf{研究背景}\ 分布式系统中的日志和指标通常预先定义且以单组件为中心，发生未知问题时，开发者难以在运行中提出新的跨组件问题，也难以关联不同位置的事件。\par
\noindent\textbf{研究问题}\ 研究问题一（RQ1）：能否动态定义监测查询，并在跨机器和跨组件事件之间保留因果关系？\par
研究问题二（RQ2）：动态因果监测能否以低开销定位软件缺陷、错误配置和性能退化等不同根因？\par
\noindent\textbf{研究方法}\ 作者结合动态插桩与先发生关系连接（happened-before join），使用户可在系统某处定义指标，却按其他组件中的因果先行事件选择、过滤和分组；原型在包含HDFS、HBase、MapReduce和YARN的Hadoop集群中评估。\par
\noindent\textbf{主要发现}\ Pivot Tracing能够跨层查询并定位软件错误、配置问题和慢硬件等多类根因，同时保持较低运行开销。它表明关系可以由明确的因果传播规则计算，而不一定逐条直接记录。\par
\noindent\textbf{贡献与意义}\ 论文为本研究的“按规则推导”提供清晰先例：关系有可检查规则和输入证据，但仍不同于原生日志事实。界面应把推导方法、覆盖范围和无法关联的事件一起呈现。\par
\clearpage
\begin{center}
{\songti\bfseries\fontsize{18.5pt}{27.75pt}\selectfont 华中科技大学设计学院硕士研究生\par}
{\songti\bfseries\fontsize{18.5pt}{27.75pt}\selectfont 文献阅读卡\par}
\end{center}
\vspace{0.25em}
\noindent\underline{\makebox[\textwidth][l]{\textbf{学号：M202570493\hfill 姓名：陈思静\hfill 序号：32}}}\par
\vspace{0.35em}
\noindent\textbf{名称}\ Dapper, a Large-Scale Distributed Systems Tracing Infrastructure\par
\noindent\textbf{作者}\ Benjamin H. Sigelman; Luiz André Barroso; Mike Burrows; Pat Stephenson; Manoj Plakal; Donald Beaver; Saul Jaspan; Chandan Shanbhag\quad 译者\par
\noindent\textbf{出处}\ 【美国】Google Technical Report dapper-2010-1，2010\quad 页\par
\vspace{0.35em}
\noindent\textbf{文献阅读心得（每篇不少于150字）：}\par
\noindent\textbf{研究背景}\ 大规模互联网服务由不同团队、语言和机器上的组件组成。开发者需要理解一次请求跨组件的路径和耗时，但全面插桩会带来开发负担、运行开销与海量数据。\par
\noindent\textbf{研究问题}\ 研究问题一（RQ1）：如何在低开销、应用透明和大规模持续部署之间取得平衡，记录端到端请求轨迹？\par
研究问题二（RQ2）：采样式分布追踪是否能支持开发和运维团队理解系统行为并构建多种诊断工具？\par
\noindent\textbf{研究方法}\ Dapper在少量公共库中插桩，以追踪标识（trace）和跨度标识（span）记录请求层级及时间信息，通过采样控制开销；论文总结系统在Google内部两年多的构建、部署、使用和衍生工具经验。\par
\noindent\textbf{主要发现}\ 公共库插桩和采样使追踪能够在超大规模生产环境持续运行，并支持性能分析、依赖发现和异常诊断。但采样和未插桩区域意味着轨迹天然不完整。\par
\noindent\textbf{贡献与意义}\ 论文建立了现代分布式追踪的追踪/跨度标识基础，并揭示观测证据的覆盖边界。对本研究而言，界面必须避免把被采样或未插桩造成的空白解释成机制没有运行。\par
\clearpage
\begin{center}
{\songti\bfseries\fontsize{18.5pt}{27.75pt}\selectfont 华中科技大学设计学院硕士研究生\par}
{\songti\bfseries\fontsize{18.5pt}{27.75pt}\selectfont 文献阅读卡\par}
\end{center}
\vspace{0.25em}
\noindent\underline{\makebox[\textwidth][l]{\textbf{学号：M202570493\hfill 姓名：陈思静\hfill 序号：33}}}\par
\vspace{0.35em}
\noindent\textbf{名称}\ A Controlled Experiment for Program Comprehension through Trace Visualization\par
\noindent\textbf{作者}\ Bas Cornelissen; Andy Zaidman; Arie van Deursen\quad 译者\par
\noindent\textbf{出处}\ 【美国】IEEE Transactions on Software Engineering，2011，37(3)：341–355，DOI: 10.1109/TSE.2010.47\quad 页\par
\vspace{0.35em}
\noindent\textbf{文献阅读心得（每篇不少于150字）：}\par
\noindent\textbf{研究背景}\ 执行轨迹可视化常被用于程序理解，但早期工具多依赖案例演示，缺少把可视化与常规开发环境直接比较的量化证据。\par
\noindent\textbf{研究问题}\ 研究问题一（RQ1）：大型执行轨迹可视化是否能提高典型程序理解任务的正确率？\par
研究问题二（RQ2）：在提供相同目标系统的情况下，可视化是否能减少完成理解任务所需的时间？\par
\noindent\textbf{研究方法}\ 作者设计8项面向代表性软件系统的程序理解任务，以Eclipse作为对照条件，以Eclipse加Extravis执行轨迹可视化作为实验条件，比较完成时间和答案正确性。\par
\noindent\textbf{主要发现}\ 使用轨迹可视化的实验组完成任务所需时间降低22\%，正确性提高43\%，两项差异均达到统计显著。研究证明结构化轨迹的价值可以用客观理解任务而非满意度评价。\par
\noindent\textbf{贡献与意义}\ 论文为本研究的三条件受控实验提供直接方法学依据：应保持任务和证据可比，测量理解正确率与时间，并把“可视化好看”与“实际理解更准确”区分开。\par
\clearpage
\begin{center}
{\songti\bfseries\fontsize{18.5pt}{27.75pt}\selectfont 华中科技大学设计学院硕士研究生\par}
{\songti\bfseries\fontsize{18.5pt}{27.75pt}\selectfont 文献阅读卡\par}
\end{center}
\vspace{0.25em}
\noindent\underline{\makebox[\textwidth][l]{\textbf{学号：M202570493\hfill 姓名：陈思静\hfill 序号：34}}}\par
\vspace{0.35em}
\noindent\textbf{名称}\ Debugging Strategies and Tactics in a Multi-Representation Software Environment\par
\noindent\textbf{作者}\ Pablo Romero; Benedict du Boulay; Richard Cox; Rudi Lutz; Sallyann Bryant\quad 译者\par
\noindent\textbf{出处}\ 【英国】International Journal of Human-Computer Studies，2007，65(12)：992–1009，DOI: 10.1016/j.ijhcs.2007.07.005\quad 页\par
\vspace{0.35em}
\noindent\textbf{文献阅读心得（每篇不少于150字）：}\par
\noindent\textbf{研究背景}\ 调试环境同时提供代码、图形和运行输出时，用户不仅要理解每种表示，还要在表示之间切换和建立对应关系。表示越多并不自动带来更好表现。\par
\noindent\textbf{研究问题}\ 研究问题一（RQ1）：程序员如何把高层调试策略与低层界面操作结合，并在多种表示之间组织搜索？\par
研究问题二（RQ2）：编程经验、图形表示知识和呈现形式如何影响调试准确率及策略适应？\par
\noindent\textbf{研究方法}\ 研究观察不同经验水平的程序员在多表示软件环境中完成调试任务，结合任务表现和交互过程分析他们查看代码、图形和输出的时间、切换方式及策略变化。\par
\noindent\textbf{主要发现}\ 成功者更系统地把代码与输出对应起来，并能随呈现条件调整策略；经验较少者更频繁地在表示间切换，形成重复且低效的浏览模式。掌握外部表示规则与编程经验同样重要。\par
\noindent\textbf{贡献与意义}\ 论文说明本研究的结构图、轨迹和上下文视图必须有明确对应关系，不能只增加表示数量。评价还应观察跨视图跳转、回退和关系核验，而不只记录最终答案。\par
\clearpage
\begin{center}
{\songti\bfseries\fontsize{18.5pt}{27.75pt}\selectfont 华中科技大学设计学院硕士研究生\par}
{\songti\bfseries\fontsize{18.5pt}{27.75pt}\selectfont 文献阅读卡\par}
\end{center}
\vspace{0.25em}
\noindent\underline{\makebox[\textwidth][l]{\textbf{学号：M202570493\hfill 姓名：陈思静\hfill 序号：35}}}\par
\vspace{0.35em}
\noindent\textbf{名称}\ Are Automated Debugging Techniques Actually Helping Programmers?\par
\noindent\textbf{作者}\ Chris Parnin; Alessandro Orso\quad 译者\par
\noindent\textbf{出处}\ 【加拿大】Proceedings of the 20th International Symposium on Software Testing and Analysis，2011：199–209，DOI: 10.1145/2001420.2001445\quad 页\par
\vspace{0.35em}
\noindent\textbf{文献阅读心得（每篇不少于150字）：}\par
\noindent\textbf{研究背景}\ 自动故障定位常以错误语句排名为目标，并假设程序员只要看到正确位置就能理解和修复问题。然而真实调试还包含依赖导航、状态检查、重跑和根因理解。\par
\noindent\textbf{研究问题}\ 研究问题一（RQ1）：现有自动调试研究关于程序员如何使用故障排名的假设是否符合实际调试行为？\par
研究问题二（RQ2）：调试技术应如何评价，才能反映其是否真正帮助开发者理解根因并完成修复？\par
\noindent\textbf{研究方法}\ 作者分析自动调试和故障定位方法的常见假设，并通过面向开发者的调试任务考察排名信息如何被检查、解释和用于后续导航，进而反思只以错误位置排名评价工具的做法。\par
\noindent\textbf{主要发现}\ 程序员不会机械地逐项检查排名，也不能仅凭孤立语句理解根因；即使正确语句排名靠前，用户仍需控制流、数据流、状态和上下文证据才能判断。\par
\noindent\textbf{贡献与意义}\ 论文确立了“自动定位准确”不等于“人能够理解并行动”的方法学原则。本研究因此应直接测量人的事件—机制映射和证据定位，而不能用自动算法指标替代人侧实验。\par
\clearpage
\begin{center}
{\songti\bfseries\fontsize{18.5pt}{27.75pt}\selectfont 华中科技大学设计学院硕士研究生\par}
{\songti\bfseries\fontsize{18.5pt}{27.75pt}\selectfont 文献阅读卡\par}
\end{center}
\vspace{0.25em}
\noindent\underline{\makebox[\textwidth][l]{\textbf{学号：M202570493\hfill 姓名：陈思静\hfill 序号：36}}}\par
\vspace{0.35em}
\noindent\textbf{名称}\ Characterizing Provenance in Visualization and Data Analysis: An Organizational Framework of Provenance Types and Purposes\par
\noindent\textbf{作者}\ Eric D. Ragan; Alex Endert; Jibonananda Sanyal; Jian Chen\quad 译者\par
\noindent\textbf{出处}\ 【美国】IEEE Transactions on Visualization and Computer Graphics，2016，22(1)：31–40，DOI: 10.1109/TVCG.2015.2467551\quad 页\par
\vspace{0.35em}
\noindent\textbf{文献阅读心得（每篇不少于150字）：}\par
\noindent\textbf{研究背景}\ 可视分析中的溯源信息（provenance）既可指数据来源、系统操作历史，也可指分析者的推理和洞见。概念混用会导致系统记录了大量历史，却不能回答这些记录服务什么任务。\par
\noindent\textbf{研究问题}\ 研究问题一（RQ1）：可视化与数据分析中的溯源信息可按哪些信息类型和使用目的组织？\par
研究问题二（RQ2）：不同溯源信息类型应如何影响捕获方式、界面设计和评价方法选择？\par
\noindent\textbf{研究方法}\ 作者综合多个领域的溯源系统和研究，对记录对象、语义层次、捕获方式及使用目的进行交叉整理，提出组织框架，并讨论类型、目的和评估之间的关系。\par
\noindent\textbf{主要发现}\ 溯源信息不仅是操作日志，还包括数据、可视状态、分析步骤、推理和洞见等不同层次；回溯、复现、协作、审计和理解等目的需要不同信息，不能由单一历史记录满足。\par
\noindent\textbf{贡献与意义}\ 论文为本研究区分运行事实、界面推导和人的解释提供概念工具，也说明证据状态必须与具体关系和使用目的绑定，而不能只给整条轨迹一个统一可信度。\par
\clearpage
\begin{center}
{\songti\bfseries\fontsize{18.5pt}{27.75pt}\selectfont 华中科技大学设计学院硕士研究生\par}
{\songti\bfseries\fontsize{18.5pt}{27.75pt}\selectfont 文献阅读卡\par}
\end{center}
\vspace{0.25em}
\noindent\underline{\makebox[\textwidth][l]{\textbf{学号：M202570493\hfill 姓名：陈思静\hfill 序号：37}}}\par
\vspace{0.35em}
\noindent\textbf{名称}\ Analytic Provenance for Sensemaking: A Research Agenda\par
\noindent\textbf{作者}\ Kai Xu; Simon Attfield; T. J. Jankun-Kelly; Ashley Wheat; Phong H. Nguyen; Nallini Selvaraj\quad 译者\par
\noindent\textbf{出处}\ 【美国】IEEE Computer Graphics and Applications，2015，35(3)：56–64，DOI: 10.1109/MCG.2015.50\quad 页\par
\vspace{0.35em}
\noindent\textbf{文献阅读心得（每篇不少于150字）：}\par
\noindent\textbf{研究背景}\ 意义建构包含信息搜索、假设生成、证据比较和框架修正。若系统只保存点击而不保存分析步骤及其关系，就难以支持反思、恢复中断或追踪不确定性。\par
\noindent\textbf{研究问题}\ 研究问题一（RQ1）：分析过程中的数据、操作、子任务、假设和推理应以怎样的层级被捕获和呈现？\par
研究问题二（RQ2）：分析溯源（analytic provenance）如何支持意义建构、协作以及不确定性和信任管理？\par
\noindent\textbf{研究方法}\ 论文汇总电气电子工程师学会可视化会议（IEEE VIS）相关研讨会和既有研究，围绕溯源信息的捕获、可视化和利用提出研究议程，并用分层模型讨论从低层交互到高层任务语义的推断。\par
\noindent\textbf{主要发现}\ 低层事件容易自动记录，高层任务和意图难以可靠推断；溯源信息可支持过程反思、探索引导和质量追踪，但推断、隐私、规模化与不确定性传播仍是开放问题。\par
\noindent\textbf{贡献与意义}\ 论文直接连接溯源信息与数据—框架意义建构，支持本研究把“证据从哪里来、怎样进入结论”作为界面主结构，并要求对自动推断的高层语义显示不确定性。\par
\clearpage
\begin{center}
{\songti\bfseries\fontsize{18.5pt}{27.75pt}\selectfont 华中科技大学设计学院硕士研究生\par}
{\songti\bfseries\fontsize{18.5pt}{27.75pt}\selectfont 文献阅读卡\par}
\end{center}
\vspace{0.25em}
\noindent\underline{\makebox[\textwidth][l]{\textbf{学号：M202570493\hfill 姓名：陈思静\hfill 序号：38}}}\par
\vspace{0.35em}
\noindent\textbf{名称}\ Capturing and Supporting the Analysis Process\par
\noindent\textbf{作者}\ Nazanin Kadivar; Victor Chen; Dustin Dunsmuir; Eric Lee; Cheryl Qian; John Dill; Christopher Shaw; Robert Woodbury\quad 译者\par
\noindent\textbf{出处}\ 【美国】Proceedings of the IEEE Symposium on Visual Analytics Science and Technology，2009：131–138，DOI: 10.1109/VAST.2009.5333020\quad 页\par
\vspace{0.35em}
\noindent\textbf{文献阅读心得（每篇不少于150字）：}\par
\noindent\textbf{研究背景}\ 分析者会反复调整参数、返回旧状态并探索替代路径，但多数可视分析工具只保存当前结果，不能让人查看、恢复或理解分析过程本身。\par
\noindent\textbf{研究问题}\ 研究问题一（RQ1）：如何捕获并可视化分析者的交互历史、分支和结果依赖？\par
研究问题二（RQ2）：可编辑、可回放的分析历史能否支持过程理解、替代路径探索和分析模式复用？\par
\noindent\textbf{研究方法}\ 作者开发CzSaw，把用户交互翻译为脚本，生成可编辑的历史视图和结果依赖图，并通过文档集合分析场景展示状态恢复、分支探索和分析流程复用。\par
\noindent\textbf{主要发现}\ 历史视图使用户能够查看分析进展和替代路径，依赖图则揭示每一步与后续结果之间的逻辑关系；脚本化记录还支持修改和重新执行分析步骤。\par
\noindent\textbf{贡献与意义}\ 论文提供了“时间历史+依赖关系”双重表示的经典先例，与本研究的轨迹和机制并列结构高度对应；它也提示两类视图应联动而非各自展示重复信息。\par
\clearpage
\begin{center}
{\songti\bfseries\fontsize{18.5pt}{27.75pt}\selectfont 华中科技大学设计学院硕士研究生\par}
{\songti\bfseries\fontsize{18.5pt}{27.75pt}\selectfont 文献阅读卡\par}
\end{center}
\vspace{0.25em}
\noindent\underline{\makebox[\textwidth][l]{\textbf{学号：M202570493\hfill 姓名：陈思静\hfill 序号：39}}}\par
\vspace{0.35em}
\noindent\textbf{名称}\ The Role of Uncertainty, Awareness, and Trust in Visual Analytics\par
\noindent\textbf{作者}\ Dominik Sacha; Hansi Senaratne; Bum Chul Kwon; Geoffrey Ellis; Daniel A. Keim\quad 译者\par
\noindent\textbf{出处}\ 【美国】IEEE Transactions on Visualization and Computer Graphics，2016，22(1)：240–249，DOI: 10.1109/TVCG.2015.2467591\quad 页\par
\vspace{0.35em}
\noindent\textbf{文献阅读心得（每篇不少于150字）：}\par
\noindent\textbf{研究背景}\ 可视分析系统从数据、模型到界面和人的解释都会引入或继承不确定性。若使用者看不到不确定性如何传播，就可能对结果形成过强或过弱的信任。\par
\noindent\textbf{研究问题}\ 研究问题一（RQ1）：不确定性如何在可视分析的数据处理、模型、可视化和人类推理之间传播？\par
研究问题二（RQ2）：不确定性意识如何影响证据搜索、验证循环和信任形成？\par
\noindent\textbf{研究方法}\ 作者整合可视分析、认知和信任研究，提出包含知识生成、探索、验证、不确定性传播与信任构建的概念模型，并通过分析场景说明各环节关系。\par
\noindent\textbf{主要发现}\ 不确定性会在系统组件中继承和累积，使用者的意识决定其是否继续寻找证据、质疑结果或采取行动；信任应建立在可检查证据和验证循环上，而非结果外观。\par
\noindent\textbf{贡献与意义}\ 论文为本研究的证据状态提供理论支撑：界面不仅要显示关系，还要让使用者意识到关系的证据强度与未知范围，从而减少由完整结构图诱发的过度确定。\par
\clearpage
\begin{center}
{\songti\bfseries\fontsize{18.5pt}{27.75pt}\selectfont 华中科技大学设计学院硕士研究生\par}
{\songti\bfseries\fontsize{18.5pt}{27.75pt}\selectfont 文献阅读卡\par}
\end{center}
\vspace{0.25em}
\noindent\underline{\makebox[\textwidth][l]{\textbf{学号：M202570493\hfill 姓名：陈思静\hfill 序号：40}}}\par
\vspace{0.35em}
\noindent\textbf{名称}\ Why Authors Don't Visualize Uncertainty\par
\noindent\textbf{作者}\ Jessica Hullman\quad 译者\par
\noindent\textbf{出处}\ 【美国】IEEE Transactions on Visualization and Computer Graphics，2020，26(1)：130–139，DOI: 10.1109/TVCG.2019.2934287\quad 页\par
\vspace{0.35em}
\noindent\textbf{文献阅读心得（每篇不少于150字）：}\par
\noindent\textbf{研究背景}\ 尽管不确定性表达方法很多，公开可视化仍经常给出确定性的单一叙事。问题不仅是缺少图形技术，还包括作者对受众、说服力和信息复杂度的假设。\par
\noindent\textbf{研究问题}\ 研究问题一（RQ1）：经常制作可视化的作者为何承认不确定性重要，却仍选择不直接呈现？\par
研究问题二（RQ2）：哪些工作实践、态度和传播目标维持了不确定性省略的规范？\par
\noindent\textbf{研究方法}\ 研究调查90名经常为他人制作可视化的作者，并访谈13名有影响力的可视化设计者；作者据此归纳实践障碍和信念矛盾，提出不确定性省略的修辞模型。\par
\noindent\textbf{主要发现}\ 作者普遍认可不确定性的价值，但担心它削弱清晰信息、增加受众负担或破坏叙事。现实中不确定性因此经常被省略，完整而确定的图形外观会掩盖推断自由度和证据边界。\par
\noindent\textbf{贡献与意义}\ 论文解释了为什么“证据状态应显示”并不会自然转化为界面实践。本研究需要验证三类状态怎样呈现才既不造成认知负担，也不让结构图以确定外观掩盖未知。\par
\restoregeometry
% 文献卡采用表格式左对齐；卡片结束后恢复中文正文的两字首行缩进。
\setlength{\parindent}{2em}
